% ======================================================================
% MNRAS paper draft (full rewrite, 2026-05-06)
%
% GPU-accelerated heterodyned nested-sampling re-analysis of the
% GW170817 bright-siren H_0 measurement of Abbott et al. (2017,
% Nature 551, 85; arXiv:1710.05835), focused on the validity of the
% post-hoc prior-reweighting step used in that work. Primary
% scientific deliverables:
%
%   1. A controlled prior-sensitivity test of the GW170817 bright-siren
%      H_0 measurement under the locked primary waveform
%      IMRPhenomXAS_NRTidalv3, including a direct-vs-reweighted
%      uniform-in-$d_L$ comparison (with effective-sample-size
%      diagnostic), a host peculiar-velocity uncertainty and centre
%      sweep, and a $d_L$-$\iota$ bimodality decomposition that
%      mechanistically explains the reweighting deficit.
%   2. A two-waveform cross-check (IMRX, TF2) bounding the residual
%      waveform systematic.
%   3. Heterodyne speedup and live-point scaling on a single A100 GPU.
%
% Numerical claims in Tables 1--4 are all generated by
% Plots/build_paper_tables.py from the canonical sample CSVs and the
% sampler.log lnZ entries; provenance is recorded in
% paper_knowledge_base/result_link_index.md.
% ======================================================================

\documentclass[fleqn,usenatbib]{mnras}

\usepackage[T1]{fontenc}
\usepackage{newtxtext,newtxmath}
\usepackage{graphicx}
\usepackage{amsmath}
\usepackage{booktabs}
\usepackage{url}
\usepackage{xspace}
\usepackage{xcolor}

\graphicspath{%
  {figures/}%
  {../results/gwtc1_phasemarg/plots/}%
}

\newcommand{\kmsmpc}{\ensuremath{\,\rm km\,s^{-1}\,Mpc^{-1}}\xspace}
\newcommand{\kms}{\ensuremath{\,\rm km\,s^{-1}}\xspace}
\newcommand{\hzero}{\ensuremath{H_{0}}\xspace}
\newcommand{\dL}{\ensuremath{d_{L}}\xspace}
\newcommand{\sigmavp}{\ensuremath{\sigma_{v_p}}\xspace}
\newcommand{\IMR}{IMRPhenomD\_NRTidalv2\xspace}
\newcommand{\IMRX}{IMRPhenomXAS\_NRTidalv3\xspace}
\newcommand{\XPHM}{IMRPhenomXPHM\xspace}
\newcommand{\TF}{TaylorF2\xspace}
\newcommand{\bjns}{\textsc{blackjax-ns}\xspace}
\newcommand{\bilby}{\textsc{bilby}\xspace}
\newcommand{\pbilby}{\textsc{parallel\_bilby}\xspace}
\newcommand{\jax}{\textsc{jax}\xspace}
\newcommand{\ripple}{\textsc{ripple}\xspace}

\title[Rapid GW170817 \hzero inference]{Rapid Hubble constant inference from GW170817 using GPU-accelerated nested sampling: prior sensitivity and the limits of post-hoc reweighting}

\author[Yang, Prathaban, Yallup \& Handley]{
Ming Yang,$^{1}$\thanks{E-mail: mhy32@cantab.ac.uk}
Metha Prathaban,$^{2,3}$
David Yallup,$^{4,2}$
and Will Handley$^{4,2}$
\\
$^{1}$Independent researcher (formerly Institute of Astronomy and St John's College, University of Cambridge), UK\\
$^{2}$Kavli Institute for Cosmology, University of Cambridge, Madingley Road, Cambridge CB3 0HA, UK\\
$^{3}$Department of Physics, University of Cambridge, JJ Thomson Avenue, Cambridge CB3 0HE, UK\\
$^{4}$Institute of Astronomy, University of Cambridge, Madingley Road, Cambridge CB3 0HA, UK
}

\date{Accepted XXX. Received YYY; in original form ZZZ}
\pubyear{2026}

\begin{document}
\label{firstpage}
\pagerange{\pageref{firstpage}--\pageref{lastpage}}
\maketitle

% ======================================================================
\begin{abstract}
The bright-siren measurement of the Hubble constant from GW170817 \citep{Abbott2017H0} assumes that switching from a volumetric to a uniform-in-$d_L$ luminosity-distance prior can be implemented by post-hoc reweighting of the baseline samples, rather than by re-running the inference under the target prior. Using a GPU-native heterodyned nested-sampling pipeline that completes the full $n_{\rm live}=5000$ analysis in $\sim 13\,\rm min$ on a single A100, we recompute the GW170817 \hzero posterior under four prior variants for the modern aligned-spin tidal waveform IMRPhenomXAS\_NRTidalv3. Switching from the volumetric to a uniform-in-$d_L$ distance prior raises the high-tail probability $P(\hzero>120\kmsmpc)$ from $0.017$ to $0.159$ when imposed during sampling and shifts the weighted-median \hzero\ from $77.6$ to $87.6\kmsmpc$, while the binned MAP stays at $70.5\kmsmpc$: both the tail and the bulk move under a change of prior that leaves the mode in place. Post-hoc reweighting of the baseline samples to the same target prior recovers only $P=0.041$ in the tail, approximately $17\,\%$ of the directly sampled shift. The three prior variants that carry an independent nested-sampling evidence agree to $\Delta\ln Z\lesssim 1.8$, so the data show at most a weak preference among the distance priors; the tail and bulk shifts are therefore properties of the prior, not a data update. Targeted mode-isolated runs reveal an $(\dL,\iota)$ bimodality whose high-\hzero, low-\dL\ branch (Mode~B; $|\ln\mathcal{B}_{\rm B/A}|<1$) the volumetric prior assigns negligible mass: this is the mechanism behind the reweighting deficit. The reweighted posterior has a lower effective sample size than the baseline, independently flagging the coverage failure. The runtime budget makes full-sample prior-sensitivity reruns the default robustness tool for bright-siren cosmology, replacing post-hoc reweighting.
\end{abstract}

\begin{keywords}
gravitational waves -- methods: data analysis -- cosmological parameters -- stars: neutron -- distance scale
\end{keywords}

% ======================================================================
\section{Introduction}
\label{sec:introduction}

The present-day expansion rate of the Universe, quantified by the Hubble constant \hzero, is the subject of one of the most persistent tensions in contemporary cosmology. Measurements anchored in the early Universe, derived from cosmic-microwave-background observations assuming flat $\Lambda$CDM, favour a low value \citep{Planck2020}, while late-Universe distance-ladder measurements based on Cepheids and Type Ia supernovae find a significantly higher value \citep{Riess2016,Riess2022}. The statistical and systematic discrepancy between these two families of measurement is now established at several $\sigma$, and its resolution -- whether through unrecognised systematics or new cosmological physics -- motivates independent probes of the expansion rate.

Gravitational-wave standard sirens are such a probe. The waveform amplitude of a compact-binary merger fixes the luminosity distance directly, without reference to an astrophysical distance ladder \citep{Schutz1986}; \citet{PalmeseMastrogiovanni2025} review the methodology and its current status. When the merger is accompanied by an electromagnetic counterpart that identifies a host galaxy, the host redshift can be combined with the gravitational-wave distance to infer \hzero from a single event \citep{HolzHughes2005}. The binary neutron-star merger GW170817 \citep{Abbott2017GW170817Discovery} is the canonical bright siren and the only confirmed example to date. Its association with the kilonova AT2017gfo in NGC~4993 yielded the first gravitational-wave \hzero measurement, $\hzero=70_{-8}^{+12}\kmsmpc$ at $68\,\%$ HPD \citep{Abbott2017H0}. That measurement has since been revisited from several directions: superluminal-jet constraints on the binary inclination \citep{Mooley2018Nature,Hotokezaka2019}, improved peculiar-velocity modelling \citep{Nicolaou2020,Mukherjee2021Velocity,HowlettDavis2020}, and an updated multi-wavelength afterglow analysis that tightens the single-event result to $\hzero=75.5_{-5.4}^{+5.3}\kmsmpc$ at $\sim 7\,\%$ precision \citep{Palmese2024GW170817H0}. A parallel line of work has argued that residual binary-viewing-angle uncertainty contributes a $\gtrsim 10\,\%$ systematic on the GW170817 \hzero \citep{SalvareseChen2024}; the prior-sensitivity question we address here sits in the same systematics-quantification programme.

For a single bright siren, the \hzero uncertainty is dominated by the luminosity-distance--inclination degeneracy: a face-on binary at larger \dL can produce the same observed strain as a more inclined binary at smaller \dL, with the two solutions mapping to different values of $\hzero \propto (v_{\rm r}-v_{p})/\dL$ for a fixed host recessional velocity $v_{\rm r}$. The resulting posterior is broad and skewed, and only weakly informative on \hzero from one event. In such a regime, the measurement is more sensitive to prior choices than the statistical uncertainty alone would suggest, because the prior shapes the tails of the posterior rather than just its peak. The original GW170817 analysis adopted a volumetric distance prior $\pi(\dL)\propto\dL^{2}$ over $[10,75]\,\rm Mpc$ and reported that switching to a uniform-in-$d_L$ prior, implemented by \emph{reweighting} posterior samples, left the bulk of the inferred \hzero\ posterior consistent with the volumetric baseline (their Extended Data Table~1; \citealp{Abbott2017H0}). Whether the same conclusion holds when the prior is imposed during sampling rather than after the fact is the question this paper addresses.

Reweighting is attractive because it avoids rerunning expensive Bayesian inference. For a posterior $p(\theta\mid d,\pi_{0})$ obtained under baseline prior $\pi_{0}$, the target posterior under an alternative prior $\pi_{1}$ can be estimated by assigning each sample a weight proportional to $\pi_{1}(\theta)/\pi_{0}(\theta)$. This procedure is statistically consistent provided the baseline samples cover the support of the target posterior. If $\pi_{1}/\pi_{0}$ is large in regions that the baseline run has sampled sparsely, the reweighted estimator has high variance and systematically low tail mass \citep[see e.g.][]{Speagle2020Dynesty}. For GW170817, the baseline volumetric prior places little mass at low \dL; a uniform-in-$d_L$ target up-weights exactly this region, and the validity of reweighting is not obvious a priori.

Testing reweighting against a direct posterior obtained by sampling under the target prior requires repeated full gravitational-wave parameter estimation. In CPU frameworks such as \bilby \citep{Ashton2019Bilby} and \pbilby, single-event binary-neutron-star analyses with phase-marginalised frequency-domain likelihood are sufficiently expensive that systematic exploration of distance, velocity, sampler, and waveform assumptions is rarely carried out in practice. Several recent developments make this constraint less severe: relative binning, also known as the heterodyned likelihood \citep{Cornish2010,Zackay2018RelativeBinning}, reduces the effective frequency grid for a long inspiral to a small number of bins; differentiable waveform libraries in \jax \citep{Bradbury2018JAX,Edwards2023Ripple} enable end-to-end GPU evaluation of the likelihood; and accelerated nested-sampling kernels, both CPU-side via normalising-flow proposals \citep{Williams2021Nessai} and GPU-native via vectorised slice sampling \citep{BlackJAX,Yallup2025BlackJAXNS,Prathaban2025GPUSpeed}, retain the evidence estimates that nested sampling provides for free while permitting large live-point counts and repeated full reruns. Related accelerated binary-neutron-star pipelines \citep{Krishna2023RelativeBinningBilby,Wong2023Jim,Wouters2024JimBNS} and forward-looking analyses \citep{HuVeitch2025} underscore the practical importance of compute-efficient parameter estimation for the third-generation detector era.

In this paper we use a GPU-native nested-sampling pipeline, built on \bjns \citep{Yallup2025BlackJAXNS,Prathaban2025GPUSpeed} and a \jax\ implementation of the heterodyned likelihood, to revisit the GW170817 bright-siren \hzero inference of \citet{Abbott2017H0}. Our focus is not a new single-event measurement of \hzero, but a controlled, multi-axis assessment of how that measurement depends on its priors, on the waveform, on the sampler hyperparameters, and on auxiliary inputs. We first validate the pipeline against GW150914 with the LVK GWTC-2.1+ production waveform \XPHM \citep{Pratten2021XPHM,Abbott2021GWTC2p1}. We then report the \IMRX\ baseline GW170817 posterior, compare it with a direct uniform-in-$d_L$ posterior and with the reweighted uniform-in-$d_L$ posterior, and use targeted mode-isolated runs to identify the $(\dL,\iota)$ bimodality that mechanistically explains the reweighting deficit. A waveform cross-check against \TF\ bounds the residual waveform systematic. Robustness sweeps over sampler hyperparameters, heterodyne-bin count, PSD source, heterodyne reference parameters, the host peculiar-velocity centre, and a companion \IMR\ prior sweep are referenced in Appendix~\ref{app:robustness}. We finish by reporting heterodyne speedup and live-point scaling.

Section~\ref{sec:method} describes the inference setup. Section~\ref{sec:validation} reports validation. Section~\ref{sec:results} contains the prior-sensitivity analysis (the central science finding), the bimodality characterisation, and the waveform cross-check. Section~\ref{sec:performance} reports runtime and live-point scaling. Section~\ref{sec:discussion} discusses implications and limitations. Section~\ref{sec:conclusions} concludes. Appendix~\ref{app:robustness} reports the inline robustness sweeps that are referenced in the body but not promoted to main figures.

% ======================================================================
\section{Method}
\label{sec:method}

\subsection{Nested sampling}

For data $d$ and source parameters $\theta$ under model $\mathcal{M}$, we target the posterior
\begin{equation}
p(\theta\mid d,\mathcal{M}) = \frac{\mathcal{L}(d\mid\theta,\mathcal{M})\,\pi(\theta\mid\mathcal{M})}{Z(d\mid\mathcal{M})},
\end{equation}
where $\mathcal{L}$ is the gravitational-wave likelihood, $\pi$ the prior, and $Z$ the Bayesian evidence. We estimate $Z$ and draw posterior samples by nested sampling \citep{Skilling2006Nested}, using the GPU-native slice sampler implemented in \bjns \citep{Yallup2025BlackJAXNS,Prathaban2025GPUSpeed}. The sampler evolves a population of live points through batched slice updates; this vectorisation keeps the GPU saturated when the per-evaluation likelihood cost is small.

All heterodyned GW170817 science runs reported here use $n_{\rm live}=5000$, $n_{\rm delete}=2500$ points removed per iteration ($\tfrac12\,n_{\rm live}$), and $n_{\rm mcmc}=8\,n_{\rm dim}=112$ MCMC steps per slice update; runs terminate when the fractional evidence increment is below $10^{-3}$. Scaling experiments vary $n_{\rm live}$ while holding all other settings fixed. The GW150914 \XPHM\ validation run uses the larger $n_{\rm live}=8000$, $n_{\rm delete}=2400$, $n_{\rm mcmc}=160$ ($16\,n_{\rm dim}$) configuration to verify that the validation conclusions are not sampler-limited at the smaller short-signal $n_{\rm dim}$. A dedicated $n_{\rm mcmc}$ convergence sweep at $\{5,10,20\}\times n_{\rm dim}$ for the 14-dimensional GW170817 problem is left to follow-up work, with the headline tail probability $P(\hzero>120\kmsmpc)$ as the natural convergence diagnostic.

\subsection{Heterodyned likelihood}

The full frequency-domain likelihood for GW170817 uses $259\,201$ bins over $20$--$2048\,\rm Hz$ at $\Delta f=1/128\,\rm Hz$ for the $128\,\rm s$ strain segment. We use a heterodyned (relative-binning) likelihood \citep{Cornish2010,Zackay2018RelativeBinning}, in which the ratio between a candidate waveform and a reference waveform is approximated by a piecewise-linear function over a small set of coarse frequency bins; we use $501$ nominal bins for GW170817 and $383$ for GW150914. The reference waveform parameters and the per-waveform effective bin count after frequency-support masking are recorded in each run's \texttt{config.json}. We compare heterodyned and unheterodyned posteriors directly in Section~\ref{sec:hetero-vs-unhet} to verify that the approximation does not bias the source-parameter posterior at the level needed for \hzero inference.

All heterodyned runs analytically marginalise over the coalescence phase, so the sampled parameter space is fourteen-dimensional for GW170817 (intrinsic masses, aligned-spin $z$-components, tidal deformabilities, inclination, luminosity distance, coalescence time, polarisation angle, sky position, and the cosmological parameters $\hzero$ and $v_{p}$), and ten-dimensional for GW150914.

\subsection{Joint Hubble-constant likelihood}

We sample the Hubble constant jointly with the gravitational-wave source parameters rather than post-processing the \dL\ marginal. The GW170817 host galaxy NGC~4993 contributes a CMB-frame recessional velocity $v_{\rm r}=3327\kms$ with $\sigma_{v_{\rm r}}=72\kms$ and a Hubble-flow peculiar-velocity estimate $\langle v_{p}\rangle=310\kms$ with $\sigmavp=150\kms$. Both enter as Gaussian likelihood factors,
\begin{equation}
\ln\mathcal{L}_{\rm tot}(\theta,\hzero,v_{p}) = \ln\mathcal{L}_{\rm GW}(d\mid\theta) + \ln\mathcal{N}\!\left(v_{\rm r}\,;\,\hzero\,\dL+v_{p},\,\sigma_{v_{\rm r}}\right) + \ln\mathcal{N}\!\left(\langle v_{p}\rangle\,;\,v_{p},\,\sigmavp\right),
\label{eq:h0likelihood}
\end{equation}
following \citet{Abbott2017H0}. Equation~\ref{eq:h0likelihood} carries no $1/N_{\rm s}(\hzero)$ GW-detection selection factor because, for every prior variant considered in this work, $\pi(\dL\mid\hzero)$ is independent of \hzero\ and the selection normalisation $N_{\rm s}=\int P_{\rm det}(\dL)\,\pi(\dL\mid\hzero)\,d\dL$ is therefore an \hzero-independent constant that cancels from the joint posterior; we return to this point with the prior definitions in Section~\ref{sec:priors}.

\subsection{Priors and prior variants}
\label{sec:priors}

The baseline prior configuration follows the bright-siren analysis of \citet{Abbott2017H0}: $\pi(\dL)\propto\dL^{2}$ over $[10,75]\,\rm Mpc$, $\pi(\hzero)\propto\hzero^{-1}$ over $[45,250]\kmsmpc$, and $\pi(v_{p})\propto 1$ over $[-1000,1000]\kms$ combined with the Gaussian $v_{p}$-likelihood term in equation~\ref{eq:h0likelihood} at $\sigmavp=150\kms$. The source-parameter priors reproduce those of the LVK GW170817 parameter estimation \citep{Abbott2019GW170817Properties}: the prior is uniform in the detector-frame component masses over $[0.5,7.7]\,M_{\odot}$ with $m_{1}\ge m_{2}$, subject to the additional chirp-mass constraint $1.184\le\mathcal{M}^{\rm det}\le 2.168\,M_{\odot}$ \citep[][Sec.~III\,D]{Abbott2019GW170817Properties}. As noted there, these bounds were chosen for technical reasons and the posterior has no support near them. The prior is full-sky, with right ascension uniform on $[0,2\pi]$ and declination distributed according to a $\cos\delta$ prior on $[-\pi/2,\pi/2]$. Every heterodyned run reported here is full-sky.

We consider three prior variants in addition to the baseline:
\begin{enumerate}
    \item \emph{Direct uniform-in-$d_L$}: $\pi(\dL)$ uniform over the same fixed $[10,75]\,\rm Mpc$ window (independent of \hzero), sampled directly.
    \item \emph{Reweighted uniform-in-$d_L$}: the baseline posterior samples reweighted by the analytic ratio of the uniform-in-$d_L$ to volumetric \dL\ priors.
    \item \emph{Enlarged peculiar-velocity uncertainty}: identical to the baseline except $\sigmavp=250\kms$.
\end{enumerate}
The direct and reweighted uniform-in-$d_L$ posteriors share the same target prior; they differ only in whether the prior is imposed during sampling or applied post-hoc. At GW170817's source redshift ($z\lesssim 0.02$), the uniform-in-$d_L$ prior is numerically equivalent within one per cent to the ``flat-in-redshift'' variant analysed by \citet{Abbott2017H0} over the same distance range; we adopt the literal uniform-in-$d_L$ form because it is what the inference pipeline imposes ($\pi(\dL)$ is a fixed density on the observable \dL, with no \hzero-dependent window). Consequently $\pi(\dL\mid\hzero)$ has no \hzero-dependence for any of the three prior variants, so the GW-detection selection normalisation $N_{\rm s}(\hzero)$ of equation~\ref{eq:h0likelihood} is constant in \hzero\ and cancels from the joint posterior, exactly as for the volumetric baseline.\footnote{We verified this property by direct evaluation of $N_{\rm s}(\hzero)=\int P_{\rm det}(\dL)\,\pi(\dL\mid\hzero)\,d\dL$ over the LVK-matched $\hzero\in[45,250]\kmsmpc$ prior with a Finn--Chernoff antenna-pattern detection model \citep{FinnChernoff1993}: $N_{\rm s}$ is \hzero-independent to machine precision for every horizon distance tested ($D_{\rm h}\in\{100,150,220\}\,\rm Mpc$, spanning the published O2 BNS range to inspiral horizon; \citealp{LVK2020Prospects}). The supporting calculation is in \texttt{mnras\_paper/test\_suite/analysis/analyze\_selection\_term.py} in the public repository.}

\subsection{Waveforms}
\label{sec:waveforms}

For GW170817 we adopt two waveforms in the main analysis. The primary is \IMRX, the most recent NR-calibrated tidal phase prescription on a modern aligned-spin base. \TF\ is included as a post-Newtonian inspiral-only family check. We additionally analysed \IMR\ as an anchor waveform for the IMRX-vs-NRTidalv2 capture-fraction comparison reported inline in Section~\ref{sec:prior}; the corresponding full posteriors are retained in the public test suite and summarised in Appendix~\ref{app:robustness}.

For GW150914 validation we use \XPHM \citep{Pratten2021XPHM}, the LVK GWTC-2.1+ production waveform for that event, which incorporates precession through a multi-scale-analysis treatment together with the dominant higher-order modes. Like-for-like comparison against the public LVK PE \citep{Abbott2021GWTC2p1,GWTC2p1_GW150914_Zenodo} therefore requires no waveform substitution.

\subsection{Posterior summaries}
\label{sec:summaries}

We summarise one-dimensional marginals by the maximum a posteriori (MAP) value and a sample-derived highest-posterior-density (HPD) interval at the requested credibility. The MAP is taken from the bin centre of the weighted histogram with maximum weighted count, on a fixed $1\kmsmpc$ grid over $[40,230]\kmsmpc$; HPDs are the shortest contiguous interval containing the requested cumulative weight, computed directly from the weighted samples. We do not use kernel-density estimates for one-dimensional summaries because the GW170817 posterior tails carry the prior-sensitivity signal of Section~\ref{sec:prior} and KDE smoothing distorts the tails. We use KDEs only for two-dimensional corner contours, where the smoothing is geometrically justified.

% ======================================================================
\section{Validation}
\label{sec:validation}

\subsection{GW150914}
\label{sec:gw150914}

We validate the pipeline on the binary-black-hole event GW150914 \citep{Abbott2016GW150914}, which has a short signal and well-established public posteriors. We analyse the public GW150914 strain data and PSDs from the LVK GWTC-2.1 Zenodo deposit \citep{GWTC2p1_GW150914_Zenodo} with \XPHM, using the same heterodyned sampler as for GW170817 with phase marginalisation; for this validation run we increase the sampler settings to $n_{\rm live}=8000$ with $n_{\rm mcmc}=160$ slice steps per update, so that the comparison is not sampler-limited at the short-signal $n_{\rm dim}$. Table~\ref{tab:gw150914} compares our recovered parameters with the LVK GWTC-2.1 public PE \citep{Abbott2021GWTC2p1,GWTC2p1_GW150914_Zenodo}; the corresponding corner overlay is shown in Figure~\ref{fig:gw150914}. The chirp mass, mass ratio, luminosity distance, and inclination all agree with the LVK reference; in particular, the recovered chirp mass agrees to ${\sim}0.2\,\%$ ($\mathcal{M}_c=30.35$ vs $30.4\,M_\odot$) and $d_{L}^{\rm med}=455\,\rm Mpc$ for our \XPHM\ run is within ${\sim}8\,\rm Mpc$ of the LVK value of $463\,\rm Mpc$. The 1$\sigma$ widths match closely: $\sigma(\mathcal{M}_c)=1.02$ vs $1.00\,M_\odot$, $\sigma(d_L)=84$ vs $85\,\rm Mpc$, $\sigma(q)=0.10$ vs $0.11$, and $\sigma(\iota)=0.41$ vs $0.35\,\rm rad$ (ratios $1.02$, $0.99$, $0.89$, $1.17$); the largest residual is on the precession-sensitive $\iota$. We use a component-mass prior uniform on $[5,100]\,M_\odot$ that encompasses the LVK GWTC-2.1 $[10,80]\,M_\odot$ range, with no posterior mass near either boundary; the validation therefore reproduces the LVK posteriors within the recovered support.

\begin{table}
  \centering
  \caption{GW150914 validation: median values of the principal source parameters for our heterodyned \XPHM\ run, with a $n_{\rm live}=5000$ cross-check. The LVK GWTC-2.1 \XPHM\ public PE \citep{Abbott2021GWTC2p1,GWTC2p1_GW150914_Zenodo} reports $\mathcal{M}_c\approx 30.4\,M_\odot$, $q\approx 0.85$, $d_L\approx 463\,\rm Mpc$, $\iota\approx 2.62\,\rm rad$ for direct comparison.}
  \label{tab:gw150914}
  \begin{tabular}{lccccc}
    \toprule
    Run & $\mathcal{M}_{c}/M_{\odot}$ & $q$ & $d_{L}/\rm Mpc$ & $\iota/\rm rad$ & $\ln Z$ \\
    \midrule
    this work, \XPHM\ ($n_{\rm live}=8000$)          & 30.35 & 0.87 & 455 & 2.61 & $260.86\pm0.06$ \\
    this work, \XPHM\ ($n_{\rm live}=5000$, cross-check) & 30.34 & 0.87 & 460 & 2.62 & $261.09\pm0.08$ \\
    \bottomrule

  \end{tabular}
\end{table}

\begin{figure}
  \centering
  \includegraphics[width=\columnwidth]{corner_GW150914_waveform_comparison.pdf}
  \caption{GW150914 validation: corner overlay of the LVK GWTC-2.1 \XPHM\ posterior \citep{Abbott2021GWTC2p1,GWTC2p1_GW150914_Zenodo} and our heterodyned \XPHM\ run at $n_{\rm live}=8000$, $n_{\rm mcmc}=160$. The two posteriors overlap throughout the recovered support; our component-mass prior $[5,100]\,M_\odot$ encompasses the LVK $[10,80]\,M_\odot$ range without adding posterior mass at the boundaries.}
  \label{fig:gw150914}
\end{figure}

\subsection{Heterodyned vs unheterodyned consistency on GW170817}
\label{sec:hetero-vs-unhet}

We next verify that the heterodyned approximation does not bias the GW170817 source-parameter posterior. We run the same sampler on the full $259\,201$-bin frequency-domain likelihood for both \IMR\ and \TF\ at $n_{\rm live}=1500$, and additionally at $n_{\rm live}\in\{500,2500\}$ for \IMR. The heterodyned and unheterodyned $\dL$, $\iota$, and $\hzero$ marginals overlay closely, and the heterodyned $\hzero^{\rm MAP}$ lies between the surrounding unheterodyned cross-checks at the level of the sampler-statistical width. We conclude that the heterodyned approximation is adequate for \hzero inference at this event's signal-to-noise ratio; the corresponding wall-clock reduction is quantified in Section~\ref{sec:speedup}.

% ======================================================================
\section{Results}
\label{sec:results}

\subsection{Prior sensitivity (the central science finding)}
\label{sec:prior}

The central scientific result of this paper is a direct comparison between the uniform-in-$d_L$ \hzero\ posterior obtained by sampling under the target prior and that obtained by reweighting the baseline posterior to the same target prior, under the locked primary waveform \IMRX. We carry out the four prior variants of Section~\ref{sec:priors} on the LVK-matched full-sky prior set and report sample-derived HPDs in Table~\ref{tab:h0priors}; the corresponding marginals are shown in Figure~\ref{fig:h0prior}.

\begin{table*}
  \centering
  \caption{\IMRX\ \hzero\ summaries on the LVK-matched full-sky prior set, $n_{\rm live}=5000$. Upper block: the four distance-prior variants of Section~\ref{sec:priors} (volumetric baseline, direct and reweighted uniform-in-$d_L$, and $\sigmavp=250\kms$). Lower block: the peculiar-velocity centre sweep $\langle v_p\rangle\in\{215,405\}\kms$ at fixed $\sigmavp=150\kms$ (the $\langle v_p\rangle=310\kms$ point is the baseline). MAP, weighted median, and sample-derived 68 per cent HPDs in $\kmsmpc$; the median is reported as a bin-noise-robust complement to the MAP for the skewed posteriors of the uniform-in-$d_L$ variants. The reweighted variant is post-hoc and has no independent $\ln Z$. All numbers generated by \texttt{Plots/build\_paper\_tables.py} from the canonical sample CSVs.}
  \label{tab:h0priors}
  \begin{tabular}{lcccccc}
    \toprule
    Variant & MAP & median & 68\% HPD & $P(\hzero>120)$ & $P(\hzero>150)$ & $\ln Z$ \\
    \midrule
    Baseline ($\pi(d_L)\propto d_L^{2}$)           & 70.5 & 77.6 & $[63.8,87.6]$ & $0.017$ & $<10^{-4}$ & $486.25\pm0.11$ \\
    Uniform-in-$d_L$, direct                       & 70.5 & 87.6 & $[64.2,103.8]$ & $0.159$ & $0.038$ & $487.30\pm0.10$ \\
    Uniform-in-$d_L$, reweighted                   & 73.5 & 82.9 & $[65.2,95.9]$ & $0.041$ & $0.000$ & (post-hoc) \\
    $\sigma_{v_p}=250\,\mathrm{km\,s^{-1}}$        & 73.5 & 78.3 & $[61.7,90.5]$ & $0.069$ & $0.015$ & $485.55\pm0.09$ \\
    \midrule
    $\langle v_p\rangle=215\,\mathrm{km\,s^{-1}}$  & 74.5 & 79.9 & $[66.1,90.4]$ & $0.021$ & $<10^{-4}$ & $485.76\pm0.12$ \\
    $\langle v_p\rangle=405\,\mathrm{km\,s^{-1}}$  & 68.5 & 76.3 & $[61.5,87.0]$ & $0.030$ & $0.000$ & $486.35\pm0.10$ \\
    \bottomrule

  \end{tabular}
\end{table*}

\begin{figure}
  \centering
  \includegraphics[width=\columnwidth]{H0_prior_sensitivity.pdf}
  \caption{Prior-sensitivity comparison for \IMRX. \emph{Left panel (a):} weighted step-histograms of \hzero\ under the four distance-prior variants of Section~\ref{sec:priors} (volumetric baseline, direct uniform-in-$d_L$, reweighted uniform-in-$d_L$, and $\sigmavp=250\kms$). The directly sampled uniform-in-$d_L$ posterior places substantially more mass at high \hzero\ than the reweighted estimate, despite both targeting the same prior. \emph{Right panel (b):} peculiar-velocity centre sweep --- \hzero\ posteriors for $\langle v_{p}\rangle\in\{215,310,405\}\kms$ at fixed $\sigmavp=150\kms$ (s18 runs). The published Abbott+2017 \hzero $=70_{-8}^{+12}\kmsmpc$ band \citep{Abbott2017H0} is shown in panel (a) as a vertical reference. Planck $\Lambda$CDM and SH0ES distance-ladder bands are shown in both panels.}
  \label{fig:h0prior}
\end{figure}

\paragraph{Direct-vs-reweighted comparison.} The direct uniform-in-$d_L$ run leaves the posterior MAP unchanged at $70.5\kmsmpc$ but broadens the 68 per cent HPD on the high side from $87.6$ to $103.8\kmsmpc$, and raises the high-tail probability
\begin{equation*}
P(\hzero>120\kmsmpc) = 0.017 \quad\text{(baseline)}\quad \to \quad 0.159 \quad\text{(direct uniform-in-}d_L\text{)},
\end{equation*}
together with $P(\hzero>150\kmsmpc)$ from $<10^{-4}$ to $0.038$. The reweighted estimate captures a small fraction of this shift: $P(\hzero>120\kmsmpc)=0.041$, equivalent to
\begin{equation}
\frac{0.041 - 0.017}{0.159 - 0.017} \approx 17\,\%
\label{eq:reweight-fraction}
\end{equation}
of the direct prior-induced change. The 68 per cent HPD upper bound under reweighting reaches $95.9\kmsmpc$, short of the directly sampled $103.8\kmsmpc$ at the same credibility. The reweighted variant is itself consistent with the uniform-in-$d_L$ posterior reported by \citet{Abbott2017H0} -- their Extended Data Table~1 gives a $68.3\,\%$ HPD $\hzero=71_{-9}^{+23}\kmsmpc$ -- so our reweighting reproduces the original analysis; the deficit we report lies in the directly sampled high-\hzero\ tail, which the reweighted procedure does not access. The three \IMRX\ prior variants that carry an independent evidence agree to $\Delta\ln Z\lesssim 1.8$; the decision-relevant pair, baseline versus direct uniform-in-$d_L$, differs by only $\Delta\ln Z\approx 1.05$. On the Jeffreys scale this is at most weak-to-substantial support, and it is within the run-to-run variability seen in the public repository for this configuration. The data therefore show at most a weak preference among the distance priors: the change in the high-\hzero\ tail across these variants is dominated by the \emph{prior}, not by a data-driven update.

\paragraph{Reweighting bias versus variance.} The reweighting deficit is bias, not high variance. A non-parametric bootstrap on the reweighted estimator (4\,000 draws at the reweighted $n_{\rm eff}=27{,}539$) gives a $95$ per cent confidence interval for $P(\hzero>120\kmsmpc)$ of $[0.037, 0.042]$ -- tight to three significant figures and excluding the directly sampled $0.159$ by ${\sim}\,100$ binomial standard errors of the reweighted estimator. The reweighted estimator has converged on its own (lower) value on this draw; running longer chains under the volumetric baseline would not close the gap. Only direct sampling under the target prior recovers the high-\hzero\ tail mass. The bootstrap and the corresponding $\hat{k}$ value below are computed by \texttt{analyze\_psis\_khat.py} in the public repository and recorded in the \texttt{paper\_diagnostics.csv} summary.

\paragraph{Effective-sample-size diagnostic.} The reweighting failure is independently flagged by the effective sample size of the reweighted posterior. We use Kish's effective sample size, $n_{\rm eff}=\left(\sum_i w_i\right)^2/\sum_i w_i^2$ \citep{Kish1965}, where $w_i$ are the (nested-sampling or importance) weights. From the same $1.73\times 10^{5}$-sample \IMRX\ baseline we measure $n_{\rm eff}=30{,}695$ (efficiency $17.8\,\%$); the directly sampled uniform-in-$d_L$ run has $n_{\rm eff}=37{,}022$ from its own $1.94\times 10^{5}$-sample draw (efficiency $19.1\,\%$), while the reweighted uniform-in-$d_L$ estimator has $n_{\rm eff}=27{,}317$ -- \emph{lower} than the baseline despite reweighting from the same draw. The reweighted-versus-baseline comparison is the coverage diagnostic; the directly sampled run is listed for context only, as it has a different sampling history. This is the conventional symptom of importance-sampling weights concentrated on a small subset of the baseline draw, the regime in which reweighting is known to be unreliable \citep{Payne2019Reweighting,Vehtari2024PSIS}. The Pareto-smoothed importance-sampling $\hat{k}$ diagnostic of \citet{Vehtari2024PSIS}, computed by fitting a generalised Pareto distribution to the upper tail of the importance ratios, gives $\hat{k}=0.68$ for the reweighted \IMRX\ draw: at the upper edge of the cautionary $0.5<\hat{k}\le 0.7$ ``high variance but consistent'' regime, just below the standard $\hat{k}>0.7$ unreliability threshold. The empirical bootstrap above shows the bias is already severe at this borderline $\hat{k}$, so for prior-sensitivity reweighting in the broad-posterior bright-siren regime we recommend reporting a bootstrap confidence interval alongside $\hat{k}$ and $n_{\rm eff}$. All three diagnostics are cheap, require no additional GPU time, and would have flagged the deficit in the original analysis from the baseline draw alone.

\paragraph{Peculiar-velocity sweep.} The $\sigmavp=250\kms$ variant shifts the \hzero\ posterior modestly: $P(\hzero>120)$ moves from $0.017$ to $0.069$. Sweeping the $v_{p}$ prior \emph{centre} over the historical literature range $\langle v_{p}\rangle\in\{215, 310, 405\}\kms$ at fixed $\sigmavp=150\kms$ (Figure~\ref{fig:h0prior}b, Table~\ref{tab:h0priors}) shifts $\hzero^{\rm MAP}$ by $6\kmsmpc$ peak-to-peak ($74.5$ at $\langle v_{p}\rangle=215\kms$ to $68.5$ at $\langle v_{p}\rangle=405\kms$) while $P(\hzero>120)$ changes by less than $0.02$. The $\langle v_{p}\rangle=310\kms$ run replicates the \IMRX\ baseline to within $0.01$ in $\ln Z$. For GW170817 with NGC~4993 identified, the dominant prior-induced uncertainty in \hzero\ is therefore the distance prior, not the peculiar-velocity prior.

\paragraph{Cross-waveform check.} Repeating the four-variant prior sweep on the legacy aligned-spin tidal model \IMR\ (NRTidalv2 calibration) gives qualitatively the same behaviour with a larger amplitude: baseline $P(\hzero>120)=0.076$; uniform-in-$d_L$ direct $P=0.281$; reweighted $P=0.195$; reweighting capture fraction $(0.195-0.076)/(0.281-0.076)\approx 58\,\%$. The newer NRTidalv3 calibration in \IMRX\ tightens the upper \hzero\ tail relative to NRTidalv2, leaving less Mode-B mass for the prior shift to redistribute and therefore both reducing the absolute prior-induced change in $P(\hzero>120)$ \emph{and} reducing the fraction of that change that reweighting can recover. The qualitative conclusion -- that reweighting systematically understates the prior-induced change in the high-\hzero\ tail -- is therefore robust across the tidal-calibration axis, and is in fact \emph{more} severe under the modern waveform than under the older one.

The shifts reported here are methodological, not cosmological: GW170817 alone is consistent with both early- and late-Universe \hzero\ measurements under every prior considered. The point is that the magnitude of this prior-induced change is a property of the inference, and reweighting alone reports a materially smaller value than direct sampling does.

\subsection{The $d_{L}$--$\iota$ bimodality}
\label{sec:bimodality}

The uniform-in-$d_L$ posterior is broad and asymmetric, with non-negligible weight at $\dL\lesssim 25\,\rm Mpc$. Before isolating the modes quantitatively, we note that the same $(\dL,\iota)$ two-peak structure is visible in \emph{both} the IMRX and the \IMR\ direct uniform-in-$d_L$ posteriors (Figure~\ref{fig:bimodality-waveform-check}): the posterior weight in the Mode-B region $\dL<30\,\rm Mpc$ is $0.325$ under \IMRX\ and $0.428$ under \IMR. The bimodality is therefore a property of the data and the uniform-in-$d_L$ prior, not the NRTidalv2 tidal-phase calibration; the reduction in Mode-B weight from \IMR\ to \IMRX\ is consistent with the tighter upper \hzero\ tail under NRTidalv3 reported in Section~\ref{sec:prior}. The mode-isolated runs that follow use \IMR/NRTidalv2 at $n_{\rm live}=5000$; the cross-waveform check above shows the bimodality is robust to the tidal calibration, so the mode-isolated Bayes factor on \IMR\ characterises the same structure that drives the \IMRX\ reweighting deficit reported in Section~\ref{sec:prior}.

To resolve whether the \IMR\ mode structure is a single broadened mode or a genuine bimodality, we run two prior-restricted variants of the direct uniform-in-$d_L$ analysis: a Mode-A run with $\dL\in[30,75]\,\rm Mpc$ and a Mode-B run with $\dL\in[10,30]\,\rm Mpc$. The $\dL=30\,\rm Mpc$ split is placed at the approximate saddle separating the two maxima in the joint $(\dL,\iota)$ posterior (Figure~\ref{fig:bimodality}). We retain the original GWTC-1 heterodyne reference, and additionally run an unrestricted $\dL\in[10,75]\,\rm Mpc$ uniform-in-$d_L$ analysis with the heterodyne reference re-anchored at Mode~B ($\dL=20\,\rm Mpc$, $\iota=2.0\,\rm rad$) to test whether the apparent Mode-B mass is a heterodyne-reference artefact.

\begin{figure}
  \centering
  \includegraphics[width=\columnwidth]{bimodality_imr_vs_imrx.pdf}
  \caption{Cross-waveform check on the $(\dL,\iota)$ bimodality: weighted joint posteriors from the unrestricted direct uniform-in-$d_L$ runs at IMRPhenomD\_NRTidalv2 (left; Mode-B-anchored heterodyne reference, s10) and \IMRX\ (right; default GWTC-1 reference, s14). Both show the same two-peak structure with the Mode-B region $\dL<30\,\rm Mpc$ carrying $0.43$ (\IMR) and $0.32$ (\IMRX) of the posterior weight: the bimodality is robust across the NRTidalv2 $\to$ NRTidalv3 calibration. The smaller IMRX Mode-B weight is consistent with the tighter upper \hzero\ tail under NRTidalv3 reported in Section~\ref{sec:prior}.}
  \label{fig:bimodality-waveform-check}
\end{figure}

Figure~\ref{fig:bimodality} shows the joint $(\dL,\iota)$ posterior under the unrestricted run, together with the per-mode \hzero\ marginals. The support is a curved arm running from $(\dL,\iota)\approx(15,1.9)\,\rm Mpc{\cdot}rad$ to $(40,2.7)\,\rm Mpc{\cdot}rad$ with two distinct local maxima. Table~\ref{tab:bimodality} reports the local evidences and per-mode \hzero\ summaries.

\begin{figure*}
  \centering
  \includegraphics[width=0.92\textwidth]{bimodality.pdf}
  \caption{The $\dL$--$\iota$ bimodality in the GW170817 uniform-in-$d_L$ posterior. Left: joint $(\dL,\iota)$ posterior from the unrestricted $\dL\in[10,75]\,\rm Mpc$ run (Mode-B-anchored heterodyne reference); the prior-restricted Mode-A and Mode-B contours from the two separate restricted-prior runs are overlaid in distinct colours, with the $\dL=30\,\rm Mpc$ prior boundary marked. Right: per-mode \hzero\ marginals as weighted step-histograms. The Planck $\Lambda$CDM and SH0ES distance-ladder bands are shown for context.}
  \label{fig:bimodality}
\end{figure*}

\begin{table*}
  \centering
  \caption{Local evidences, MAP $\hzero$ values, and sample-derived 68 per cent HPDs for the $\dL$-restricted uniform-in-$d_L$ runs. Upper block: seed=0 (primary); lower block: independent seed=1 replication (s18 runs). The Mode-B/Mode-A Bayes factor includes a $\ln(20/45)$ correction for the difference in restricted-prior $\dL$ volumes.}
  \label{tab:bimodality}
  \begin{tabular}{lccccc}
    \toprule
    Variant & $\dL$ range (Mpc) & $\hzero^{\rm MAP}\,(\kmsmpc)$ & 68\% HPD & $P(\hzero>120)$ & $\ln Z$ \\
    \midrule
    Mode A         & $[30,75]$ & 74.5 & $[66.7,88.8]$ & $<10^{-4}$ & $486.80\pm0.10$ \\
    Mode B         & $[10,30]$ & 109.5 & $[99.5,151.8]$ & $0.638$ & $486.95\pm0.09$ \\
    Unrestricted   & $[10,75]$ & 73.5 & $[62.6,116.5]$ & $0.281$ & $486.48\pm0.09$ \\
    \midrule
    Mode A (s=1)   & $[30,75]$ & 72.5 & $[66.1,88.6]$ & $<10^{-4}$ & $486.71\pm0.09$ \\
    Mode B (s=1)   & $[10,30]$ & 110.5 & $[99.0,154.5]$ & $0.646$ & $487.62\pm0.10$ \\
    Unrestr. (s=1) & $[10,75]$ & 73.5 & $[62.0,121.1]$ & $0.311$ & $487.52\pm0.09$ \\
    \bottomrule

  \end{tabular}
\end{table*}

The Bayes factor between Mode~B and Mode~A, after correcting for the volume difference between the restricted priors, is
\begin{equation*}
\ln \mathcal{B}_{\rm B/A} = (\ln Z_{\rm B}-\ln Z_{\rm A}) + \ln(20/45) = +0.15 - 0.81 = -0.66 \quad (\text{seed=0}).
\end{equation*}
An independent second-seed replication (Table~\ref{tab:bimodality}, lower block) gives $+0.91 - 0.81 = +0.10$ (seed=1). The $\ln(20/45)$ term is the log-ratio of the restricted-prior $\dL$ widths; it is the exact volume correction because the prior is uniform in $\dL$ by construction. Both seeds satisfy $|\ln\mathcal{B}_{\rm B/A}|<1$, placing the result in the ``not worth more than a mention'' range of the Jeffreys scale. We caution that the run-to-run $\ln Z$ scatter is larger than the nominal $\pm 0.1$ per-run uncertainty: the unrestricted $\ln Z$ differs by $1.04$ between the two seeds, so the $\pm 0.1$ figure should be read as a within-run statistical error, not a run-to-run reproducibility bound. The conclusion that survives this scatter is the sign-independent one -- $|\ln\mathcal{B}_{\rm B/A}|<1$ in both seeds -- so Mode~B is neither significantly favoured nor disfavoured regardless of seed. Mode~B contributes a non-negligible share of the joint posterior under any prior that gives the low-\dL\ region appreciable weight. The Mode-B-anchored unrestricted run recovers $P(\hzero>120)=0.281$ -- statistically indistinguishable from the GWTC-1-anchored \IMR\ direct uniform-in-$d_L$ run reported in Section~\ref{sec:prior} -- so the heterodyne-reference choice does not bias the Mode-A/Mode-B weight ratio.

This is the mechanism behind the reweighting capture fraction of equation~\ref{eq:reweight-fraction}. Mode~B is the high-\hzero, small-\dL\ branch of the distance--inclination degeneracy. The baseline volumetric prior $\pi(\dL)\propto\dL^{2}$ assigns Mode~B a tiny prior mass (the $[10,30]\,\rm Mpc$ slab carries $(30^3-10^3)/(75^3-10^3)\approx 6\,\%$ of the total $[10,75]\,\rm Mpc$ volumetric prior), and the corresponding region of the baseline posterior is sparsely sampled. Reweighting cannot reconstruct what is not there. Direct sampling under the uniform-in-$d_L$ prior populates Mode~B at sampler resolution, and it is precisely the missing Mode-B mass that opens the gap between the directly sampled and reweighted estimates of $P(\hzero>120)$. Posterior-coverage diagnostics that compare the baseline samples against the target prior -- such as the $n_{\rm eff}$ measurement reported in Section~\ref{sec:prior} -- should therefore be standard practice before reporting a reweighted bright-siren \hzero\ summary.

\subsection{Cross-waveform Hubble-constant posterior}
\label{sec:cross-waveform}

Figure~\ref{fig:waveform-h0} shows the GW170817 \hzero\ posterior under the LVK-matched prior (Section~\ref{sec:priors}) for \IMRX\ and \TF, together with the published Abbott+2017 GW170817 \hzero\ band as a vertical reference. All curves are weighted step-histograms; sample-derived HPDs are reported in Table~\ref{tab:waveform-h0}. Figure~\ref{fig:waveform-corner} shows the corresponding $(\mathcal{M}_{c},q,\chi_{\rm eff},\dL,\iota)$ corner overlay against the public GWTC-1 IMRPhenomPv2\_NRTidal posterior \citep{Abbott2019GWTC1,LVK_GW170817_DataRelease}.

\IMRX, the locked primary, has $\hzero^{\rm MAP}=70.5\kmsmpc$, 68 per cent HPD $[63.8,87.6]\kmsmpc$ and 95 per cent HPD $[59.4,111.3]\kmsmpc$; this is the baseline run of Section~\ref{sec:prior} (Table~\ref{tab:h0priors}), reused here as the reference waveform for the cross-waveform comparison. \TF, a post-Newtonian inspiral-only family check on the same prior, has $\hzero^{\rm MAP}=68.5\kmsmpc$ -- ${\sim}2\kmsmpc$ lower, the small offset expected for a lower-PN, point-particle inspiral family -- and 68 per cent HPD $[61.4,89.3]\kmsmpc$, overlapping \IMRX\ throughout. The published Abbott+2017 result $\hzero=70_{-8}^{+12}\kmsmpc$ \citep{Abbott2017H0} is consistent with both. \TF\ carries the heavier upper tail (95 per cent upper bound $125.6$ versus $111.3\kmsmpc$ for \IMRX, and $P(\hzero>120)=0.065$ versus $0.017$), which we attribute to the more recent NR-calibrated tides in \IMRX\ damping the high-\hzero, small-\dL\ branch of the degeneracy. The IMRX-vs-IMR capture-fraction comparison reported in Section~\ref{sec:prior} indicates that the same tide-driven tail compression accounts for the smaller absolute prior-induced shift in $P(\hzero>120)$ under \IMRX.

\begin{figure}
  \centering
  \includegraphics[width=\columnwidth]{H0_waveform_comparison.pdf}
  \caption{GW170817 \hzero\ posterior for \IMRX\ (primary) and \TF\ (family check) on the LVK-matched prior. Vertical markers are 68 per cent and 95 per cent HPD bounds derived from the weighted samples (no KDE smoothing); sample-derived HPDs are reported in Table~\ref{tab:waveform-h0}. The Planck $\Lambda$CDM and SH0ES distance-ladder values are shown for context. The published Abbott+2017 GW170817 \hzero $=70_{-8}^{+12}\kmsmpc$ band \citep{Abbott2017H0} is shown as a vertical reference.}
  \label{fig:waveform-h0}
\end{figure}

\begin{figure*}
  \centering
  \includegraphics[width=0.92\textwidth]{corner_GW170817_waveform_comparison.pdf}
  \caption{GW170817 source-parameter corner overlay for \IMRX\ and \TF\ on the LVK-matched prior, with the public GWTC-1 IMRPhenomPv2\_NRTidal posterior \citep{Abbott2019GWTC1,LVK_GW170817_DataRelease} included as a literature reference. Two-dimensional contours are KDE-smoothed; one-dimensional marginals are weighted step-histograms.}
  \label{fig:waveform-corner}
\end{figure*}

\begin{table}
  \centering
  \caption{GW170817 \hzero\ summaries for the two-waveform main set on the LVK-matched prior, $n_{\rm live}=5000$. MAP and sample-derived HPDs in $\kmsmpc$. The \IMRX\ row is the baseline run of Table~\ref{tab:h0priors}.}
  \label{tab:waveform-h0}
  \begin{tabular}{lccccc}
    \toprule
    Waveform & MAP & 68\% HPD & 95\% HPD & $P(\hzero>120)$ & $\ln Z$ \\
    \midrule
    \IMRX\ (primary)             & 70.5 & $[63.8,87.6]$ & $[59.4,111.3]$ & $0.017$ & $486.25\pm0.11$ \\
    \TF\ (family check)          & 68.5 & $[61.4,89.3]$ & $[56.9,125.6]$ & $0.065$ & $486.90\pm0.10$ \\
    \bottomrule
  \end{tabular}
\end{table}

% ======================================================================
\section{Performance and runtime studies}
\label{sec:performance}

\subsection{Wall-clock runtime}
\label{sec:wallclock}

The primary heterodyned GW170817 \hzero\ analysis with \IMRX\ and $n_{\rm live}=5000$ on the LVK-matched prior set completes in $\approx 13$\,min on a single NVIDIA A100 GPU; the corresponding \TF\ run takes $\approx 4$\,min. The full four-variant prior-sensitivity suite (baseline, direct uniform-in-$d_L$, reweighted post-hoc, $\sigmavp=250\kms$) for \IMRX\ on the LVK-matched prior set fits inside an hour. The GW150914 \XPHM\ validation run at the larger $n_{\rm live}=8000$, $n_{\rm mcmc}=160$ settings takes $\approx 5\,\rm h$ on the same hardware -- the price for both the precessing waveform and the converged short-signal sampler configuration that yields the LVK reproduction of Section~\ref{sec:gw150914}.

\subsection{Heterodyne speedup and live-point scaling}
\label{sec:speedup}

The heterodyned likelihood is the enabling approximation. Figure~\ref{fig:scaling} compares heterodyned and unheterodyned \IMR\ wall-clock time on GW170817 across $n_{\rm live}\in\{500,1000,2500,5000,10\,000,20\,000,50\,000,100\,000\}$ for the heterodyned sampler and $n_{\rm live}\in\{500,1500,2500\}$ for the unheterodyned reference, plus the \TF\ unheterodyned point at $n_{\rm live}=1500$. At matched $n_{\rm live}$ the speedup factors are
\begin{equation*}
\text{speedup}(n_{\rm live}) = \begin{cases} 31.1\times & n_{\rm live}=500, \\ 50.7\times & n_{\rm live}=1500, \\ 67.7\times & n_{\rm live}=2500. \end{cases}
\end{equation*}
The relative-binning approximation accounts for the bulk of the wall-clock reduction; the GPU saturation provided by \bjns\ is layered on top.

\begin{figure*}
  \centering
  \includegraphics[width=0.92\textwidth]{scaling_study_full.pdf}
  \caption{GW170817 wall-clock scaling. Heterodyned and unheterodyned wall-clock as a function of $n_{\rm live}$ for the \IMR\ family on the LVK-matched full-sky prior set, with the \TF\ unheterodyned point at $n_{\rm live}=1500$ for context. The heterodyned series has log--log gradient $\simeq 0.60$ over the plotted range; matched unheterodyned points are $31\times$/$51\times$/$68\times$ slower at $n_{\rm live}=500/1500/2500$. The shallow bend at $n_{\rm live}=20\,000$ is a termination-tolerance artefact and is documented in the public repository.}
  \label{fig:scaling}
\end{figure*}

The heterodyned $\hzero^{\rm MAP}$ is stable across $n_{\rm live}\geq 2500$ to within $1\kmsmpc$, and $\ln Z$ stabilises within $\pm 0.2$ over the same range. The $n_{\rm live}=5000$ configuration adopted for the science runs is therefore in the converged regime. Live-point counts up to $10^{5}$ remain accessible on a single GPU; the $n_{\rm live}=10^{5}$ heterodyned run completes in roughly four hours.

\subsection{Scope of the performance claim}

The runtime numbers above are single-A100 wall-clock on the current pipeline: the full $n_{\rm live}=5000$ \hzero\ analysis completes in under a quarter of an hour, the multi-variant prior-sensitivity suite fits inside an hour, and the live-point convergence study extends to $n_{\rm live}=10^{5}$ within an overnight budget. This budget is what makes the multi-axis robustness study of this paper scientifically routine. A like-for-like \IMRX\ \pbilby\ benchmark on matched priors and live-point counts would calibrate the speedup against a production CPU pipeline and is the appropriate cross-check for a follow-up.

% ======================================================================
\section{Discussion}
\label{sec:discussion}

\subsection{Implications for bright-siren cosmology}

The main methodological message is that posterior reweighting, although statistically consistent in the limit of infinite baseline samples, can substantially understate the prior sensitivity of bright-siren \hzero\ measurements in practice. For GW170817 under \IMRX, switching from a volumetric distance prior to a uniform-in-$d_L$ prior changes $P(\hzero>120\kmsmpc)$ by a factor of $\approx 9$ under direct sampling, while reweighting captures only $\approx 17\,\%$ of that shift (equation~\ref{eq:reweight-fraction}). The corresponding capture fraction on \IMR\ -- where the high-\hzero\ tail is heavier -- rises to $\approx 58\,\%$. The deficit is a problem of effective sample coverage rather than of statistical methodology: the baseline volumetric prior assigns Mode~B a small prior mass and reweighting cannot reconstruct what was not sampled. The reweighted posterior's effective sample size, lower than the baseline's despite reweighting from the same draw, independently flags the coverage failure and is the diagnostic we recommend as a default check before any reweighted bright-siren \hzero\ summary is reported.

For a single event the cosmological interpretation is robust to this under-estimation: both early- and late-Universe \hzero\ values lie within the GW170817 68 per cent HPD under any of the priors considered here. However, the scale of the effect matters for forecasts \citep{Chen2018Forecast} and combined analyses. When bright-siren posteriors are combined over multiple events, or when tail probabilities are compared against external priors from cosmological surveys, the prior-dependence of each single-event posterior propagates through. An analysis that quantifies that dependence through reweighting alone will systematically underestimate the prior contribution.

\subsection{Where reweighting is sufficient}

Reweighting is adequate when the target and baseline priors have similar support in the regions that carry most of the posterior mass. The $\sigmavp=250\kms$ variant in our analysis is a case of this kind: the change modifies the host-velocity likelihood smoothly over the scale of the velocity posterior, the baseline samples cover the affected region well, and the reweighted and directly sampled $\sigmavp$ posteriors agree to well within the statistical width. The uniform-in-$d_L$ case fails the same coverage test, because the baseline volumetric and target uniform-in-$d_L$ priors differ most at the edges of the \dL\ range. The mode-isolated runs of Section~\ref{sec:bimodality} make this concrete: Mode~B has $|\ln \mathcal{B}_{\rm B/A}|<1$ relative to Mode~A (confirmed over two independent seeds) under the uniform-in-$d_L$ prior, but its baseline-volumetric-prior mass is only $\approx 7\,\%$ of Mode~A's, so reweighted samples are blind to it.

\subsection{Implications for future bright sirens}

The practical consequence is that direct prior sampling should be used for bright-siren \hzero\ analyses whenever the target prior differs non-trivially in support from the baseline, rather than being used only as a cross-check. The runtime budget on a single A100 GPU now permits such reruns routinely (Section~\ref{sec:performance}); the live-point and sampler-robustness studies (Appendix~\ref{app:robustness}) show that the science-quality configuration sits in the converged regime, so the marginal cost of such a rerun is the marginal cost of one ${\sim}15\,\rm min$ run. Live-point counts up to $10^{5}$ are accessible on the same hardware, which means that these studies can be combined with population-level analyses without a substantial increase in end-to-end wall-clock time. This is directly relevant to the third-generation detector era \citep{HuVeitch2025}, where the detection rate of compact-binary mergers will make repeated, controlled posterior reanalyses the norm rather than the exception.

\subsection{Scope and natural extensions}

Two items sit outside the scope of the present analysis and mark the natural extensions. First, no waveform in our \jax\ inventory simultaneously includes precession and tides, so the locked primary is the aligned-spin tidal \IMRX; a tides-with-precession waveform in \jax\ would extend the cross-waveform check. Second, the $\dL$--$\iota$ bimodality is characterised here at the level of local evidences and mode-isolated marginals; propagating Mode~B into a population analysis, or combining it with an electromagnetic likelihood for the host inclination, is the natural next step for bright-siren cosmology built on this finding.

% ======================================================================
\section{Conclusions}
\label{sec:conclusions}

We have used a GPU-native nested-sampling pipeline to revisit the GW170817 bright-siren \hzero\ measurement of \citet{Abbott2017H0}. Switching the distance prior from volumetric to uniform-in-$d_L$ by direct sampling raises $P(\hzero>120\kmsmpc)$ from $0.017$ to $0.159$ and moves the weighted-median \hzero\ from $77.6$ to $87.6\kmsmpc$, while the binned MAP stays at $70.5\kmsmpc$; post-hoc reweighting captures only $0.041$ in the tail, $\approx 17\,\%$ of the directly sampled shift. The three \IMRX\ prior variants that carry an independent evidence agree to $\Delta\ln Z\lesssim 1.8$, so the tail and bulk shifts are properties of the prior, not data updates. The mechanism is the $\dL$--$\iota$ bimodality: a high-\hzero, small-\dL\ branch (Mode~B; $|\ln\mathcal{B}_{\rm B/A}|<1$, confirmed in two independent seeds) carries appreciable likelihood but negligible volumetric-prior mass, and reweighting cannot recover what was not sampled. The reweighted estimator's effective sample size, lower than the baseline's, independently flags the coverage failure. The runtime budget makes direct prior-sensitivity reruns the right default for bright-siren cosmology, with multi-axis robustness studies now routine on a single A100 GPU.

% ======================================================================
\section*{Acknowledgements}

This work was supported by the research environment and infrastructure of the Handley Lab at the University of Cambridge. MP is supported by the Harding Distinguished Postgraduate Scholars Programme (HDPSP). This material is based upon work supported by the Google Cloud research credits program, with the award GCP397499138. We acknowledge the LIGO--Virgo--KAGRA Collaboration for the public strain data and reference posteriors used here \citep{Abbott2019GWTC1,Abbott2021GWTC2p1,LVK_GW170817_DataRelease,LVK_H0_DataRelease,GWTC2p1_GW150914_Zenodo}.

% ======================================================================
\section*{Data Availability}

The GW170817 strain data, host-galaxy peculiar-velocity inputs, and GWTC-1 reference posteriors used here are from the LIGO--Virgo Collaboration data releases LIGO P1800061 and P1700296 \citep{LVK_GW170817_DataRelease,LVK_H0_DataRelease}. The GW150914 strain data and the GWTC-2.1 \XPHM\ reference parameter-estimation posterior are from the LVK GWTC-2.1 Zenodo deposit \citep{GWTC2p1_GW150914_Zenodo}. The analysis scripts, derived summary tables, run catalogue, and figure- and table-generation code needed to reproduce the results in this paper can be obtained from \citet{Yang2026DataRelease}. The nested-sampling chains themselves are not redistributed in that repository (each individual chain is ${\sim}100\,\rm MB$ and the full set runs to several GB); they can be regenerated from the public strain data using the \bjns\ sampler \citep{Yallup2025BlackJAXNS} and the custom GW-likelihood kernel \citep{Prathaban2025GPUSpeed}, the latter of which also contains instructions and example files for running GW analyses.

% ======================================================================
\bibliographystyle{mnras}
\bibliography{references}

% ======================================================================
\appendix

\section{Inline robustness sweeps}
\label{app:robustness}

We complete the multi-axis robustness picture with sweeps held in this appendix; the central science finding of Section~\ref{sec:prior} survives every one of them.

\begin{itemize}
    \item \emph{Sampler hyperparameters.} On the \IMR\ LVK-matched baseline at $n_{\rm live}=5000$, varying $n_{\rm delete}/n_{\rm live}\in\{0.10,0.25,0.50,0.75\}$ moves $\hzero^{\rm MAP}$ by $\le 1.2\kmsmpc$ and $\ln Z$ by $\le 0.6$; varying the heterodyne-bin count $n_{\rm bins}\in\{251,501,1001\}$ moves $\hzero^{\rm MAP}$ by $\le 0.7\kmsmpc$, with the largest pairwise Wasserstein-1 distance between the three \hzero\ posteriors below $2\kmsmpc$.
    \item \emph{PSD source.} \TF\ runs at $n_{\rm live}=5000$ on the baseline prior with GWTC-1, \texttt{kazewong}, and \bilby\ reference PSDs span $\hzero^{\rm MAP}=76.4$--$77.7\kmsmpc$ and $P(\hzero>120)=0.043$--$0.065$, again with maximum pairwise W$_{1}<2\kmsmpc$.
    \item \emph{Heterodyne reference.} Replacing the GWTC-1 reference parameters with the maximum-likelihood point of an initial low-cost optimisation produces \IMR\ posteriors that are indistinguishable through the bulk of the distribution and depart only in the high-\hzero\ tail.
    \item \emph{Peculiar-velocity centre.} Sweeping $\langle v_{p}\rangle\in\{215, 310, 405\}\kms$ at fixed $\sigmavp=150\kms$ on the \IMRX\ baseline (Figure~\ref{fig:h0prior}b, Table~\ref{tab:h0priors}) shifts $\hzero^{\rm MAP}$ by $6\kmsmpc$ peak-to-peak across the full range ($74.5$ at $\langle v_{p}\rangle=215\kms$ to $68.5$ at $\langle v_{p}\rangle=405\kms$), while $P(\hzero>120)$ changes by less than $0.02$. The $\langle v_{p}\rangle=310\kms$ run reproduces the \IMRX\ baseline to within $0.01$ in $\ln Z$, confirming run reproducibility.
    \item \emph{\IMR\ companion full sweep.} The four-variant prior sensitivity sweep on \IMR\ gives baseline $\hzero^{\rm MAP}=71.5\kmsmpc$, $P(\hzero>120)=0.076$; uniform-in-$d_L$ direct $\hzero^{\rm MAP}=73.5\kmsmpc$, $P=0.281$; reweighted $\hzero^{\rm MAP}=71.5\kmsmpc$, $P=0.195$; $\sigmavp=250\kms$ $\hzero^{\rm MAP}=70.5\kmsmpc$, $P=0.067$. The reweighting-capture fraction is $(0.195-0.076)/(0.281-0.076)\approx 58\,\%$, against $\approx 17\,\%$ for \IMRX\ (equation~\ref{eq:reweight-fraction}).
\end{itemize}
All five axes therefore produce \hzero\ shifts at least an order of magnitude smaller than the prior-induced shift documented in Section~\ref{sec:prior}.

\label{lastpage}
\end{document}
